% !PRINTED 680
\weberpage{680}
%<<CONT
the exponents $\xi_i$ remain restricted to non-negative proper fractional values, and from this there results, by the definition of the conjugate logarithms, an upper limit for $\varepsilon$ and its conjugates.

But by Theorem 2., \S.\,185 there is in $\Omega$ only a finite number of integers, and hence all the more only a finite number of units which, in absolute value, lie together with all their conjugates below a finite limit.

If we now call a unit whose exponents lie within the limits 0 and 1, with inclusion of the lower and with exclusion of the upper limit, a unit \gesperrt{reduced with respect to the system $\varepsilon_1, \varepsilon_2, \ldots, \varepsilon_{\nu-1}$}, then we have the theorem:

\gesperrt{There is only a finite number of units in $\Omega$ which are reduced with respect to an independent system $\varepsilon_1, \varepsilon_2, \ldots, \varepsilon_{\nu-1}$.}

2) The units reproduce themselves by multiplication and division. If
\[
\begin{array}{l}
\xi_1', \xi_2', \ldots, \xi_{\nu-1}' \\
\xi_1'', \xi_2'', \ldots, \xi_{\nu-1}''
\end{array}
\]
are the exponents of two units $\varepsilon', \varepsilon''$, then the sums and the differences
\[
\xi_1' \pm \xi_1'', \xi_2' \pm \xi_2'', \ldots, \xi_{\nu-1}' \pm \xi_{\nu-1}''
\]
are the exponents of the units $\varepsilon' \varepsilon''$ and $\varepsilon' : \varepsilon''$.

This follows immediately from the theorem that the logarithm of a product or of a quotient is equal to the sum or to the difference of the logarithms of the two parts.

3) The units $\varepsilon_1, \varepsilon_2, \ldots$ themselves have the exponents $1, 0, \ldots, 0; 0, 1, \ldots, 0; \ldots$, and from this it follows by 2) that every system of integers, put for $\xi_1, \xi_2, \ldots, \xi_{\nu-1}$, gives the system of exponents of a unit which can be formed by multiplication of powers of the $\varepsilon_1, \varepsilon_2, \ldots$.

4) From 2) and 3) it results that a system of exponents $\xi_1, \xi_2, \ldots, \xi_{\nu-1}$ of a unit remains a system of exponents when all its elements are multiplied by one and the same integer, and that it still retains this character even when its elements $\xi_i$ are increased or diminished by arbitrary integers. If therefore we use the symbol $(x)$ in the sense of \S.\,182, namely that it denotes the ex-
%CONT>>
